%=======================================================================
%  Research Overview Poster  --  Niels Bohr Quantum Summer School 2026
%  Jaehun Han
%
%  A1 PORTRAIT (594 mm x 841 mm), 3-column layout, beamerposter.
%
%  HOW TO USE ON OVERLEAF
%  ----------------------
%  1. New Project -> Blank Project, paste this file as main.tex.
%  2. Set the compiler to  pdfLaTeX  (Menu -> Settings -> Compiler).
%  3. Drop your figure PDFs/PNGs into a  figures/  folder and replace the
%     three \includegraphics paths below (search for "REPLACE").
%     The placeholders compile out-of-the-box with NO external files.
%
%  Fonts/packages used are all standard on Overleaf; nothing extra needed.
%=======================================================================
\documentclass[final]{beamer}

% ---- Poster geometry: A1 portrait -------------------------------------
\usepackage[size=a1,orientation=portrait,scale=1.32]{beamerposter}

% ---- Basic packages ----------------------------------------------------
\usepackage[utf8]{inputenc}
\usepackage[T1]{fontenc}
\usepackage{lmodern}
\usepackage{amsmath,amssymb}
\usepackage{graphicx}
\usepackage{tikz}
\usetikzlibrary{arrows.meta,positioning}
\usepackage{qrcode}          % email QR code (comment out if unavailable)
\usepackage{ragged2e}

%=======================================================================
%  COLOUR PALETTE  --  modern, calm, academic
%=======================================================================
\definecolor{ink}{HTML}{1A2233}        % near-black text
\definecolor{accentA}{HTML}{2A6F97}    % MBQC       (blue)
\definecolor{accentB}{HTML}{9C6644}    % Surface    (warm brown)
\definecolor{accentC}{HTML}{4C956C}    % Distillation (green)
\definecolor{band}{HTML}{22303C}       % header band
\definecolor{softbg}{HTML}{F4F6F8}     % block body background
\definecolor{rule}{HTML}{C9D3DB}

%=======================================================================
%  beamer theme tweaks  --  strip decoration, big clean blocks
%=======================================================================
\setbeamercolor{block title}{fg=white,bg=ink}
\setbeamercolor{block body}{fg=ink,bg=softbg}
\setbeamercolor{headline}{fg=white,bg=band}
\setbeamerfont{block title}{series=\bfseries,size=\large}
\setbeamertemplate{navigation symbols}{}
\setbeamertemplate{caption}[numbered]
\setbeamercolor{caption name}{fg=ink}

% Block spacing / rounded look
\setbeamertemplate{block begin}{
  \vskip.4ex
  \begin{beamercolorbox}[colsep*=.75ex,rounded=true,leftskip=1.2ex,rightskip=1.2ex]{block title}
    \usebeamerfont{block title}\insertblocktitle
  \end{beamercolorbox}
  \vskip-.2ex
  \begin{beamercolorbox}[colsep*=.75ex,rounded=true,leftskip=1.2ex,rightskip=1.2ex,vmode]{block body}
    \vskip.5ex
}
\setbeamertemplate{block end}{
    \vskip.6ex
  \end{beamercolorbox}
  \vskip1.6ex
}

% Compact, readable itemize
\setbeamertemplate{itemize item}{\raise0.5ex\hbox{\footnotesize$\bullet$}}
\setbeamertemplate{itemize subitem}{\raise0.4ex\hbox{\scriptsize$\circ$}}
\setlength{\leftmargini}{1.6em}

% ---- Convenience macros ------------------------------------------------
\newcommand{\rq}[1]{{\bfseries\color{ink}Research question.\ }#1}
\newcommand{\seclabel}[2]{% coloured pill above each column title
  \begin{beamercolorbox}[sep=0.4ex,center]{normal text}
  \colorbox{#1}{\parbox{0.9\linewidth}{\centering\color{white}\bfseries\Large #2}}
  \end{beamercolorbox}\vskip0.6ex}

% Uniform figure placeholder box (used until you drop in real figures)
\newcommand{\figplaceholder}[2]{% #1 = accent colour, #2 = caption text
  \begin{tikzpicture}
    \node[draw=#1,line width=2pt,rounded corners=6pt,fill=white,
          minimum width=0.95\linewidth,minimum height=9cm,
          align=center,text=#1]
      {\bfseries FIGURE PLACEHOLDER\\[2mm]\normalfont\color{ink}\small #2};
  \end{tikzpicture}}

%=======================================================================
%  TITLE / HEADER
%=======================================================================
\title{Three Perspectives on Efficient and Reliable Quantum Computation}
\author{Jaehun Han}
\institute{KAIST Graduate School of Quantum Science and Technology \quad$\cdot$\quad QIT Lab \quad$\cdot$\quad ETRI Quantum Computing Lab}

% Custom headline (title band with QR + affiliation)
\setbeamertemplate{headline}{
  \begin{beamercolorbox}[wd=\paperwidth,sep=0pt]{headline}
    \vskip1.2cm
    \begin{columns}[c]
      \begin{column}{0.14\paperwidth}\centering
        % --- QR CODE: encodes email; replace text if you prefer a URL ---
        \begin{beamercolorbox}[sep=0.6ex,rounded=true,center]{normal text}
          \colorbox{white}{\qrcode[height=3.6cm]{mailto:pjhhan.research@gmail.com}}
        \end{beamercolorbox}
        {\color{white}\footnotesize\vskip2mm scan to contact}
      \end{column}
      \begin{column}{0.72\paperwidth}\centering
        {\color{white}\usebeamerfont{title}\bfseries\huge \inserttitle\par}
        \vskip6mm
        {\color{white}\Large A research overview across three ongoing directions in quantum computing\par}
        \vskip8mm
        {\color{white}\LARGE\bfseries \insertauthor\par}
        \vskip3mm
        {\color{white}\large \insertinstitute\par}
      \end{column}
      \begin{column}{0.14\paperwidth}\centering
        {\color{white}\footnotesize Niels Bohr\\Quantum\\Summer School\\Odense, 2026}
      \end{column}
    \end{columns}
    \vskip1.2cm
  \end{beamercolorbox}
}

%=======================================================================
\begin{document}
\begin{frame}[t]

% ---- Unifying statement banner ----------------------------------------
\begin{beamercolorbox}[wd=\textwidth,sep=1.2ex,rounded=true,center]{block body}
  \large\itshape
  Across three projects, I study how \textbf{structural information} in quantum
  computation can be exploited to \textbf{reduce resources}, \textbf{extract hidden
  error information}, and \textbf{improve noisy quantum states.}
\end{beamercolorbox}
\vskip1.4ex

\begin{columns}[t]

%=====================================================================
%  COLUMN 1 -- MBQC RESOURCE OPTIMIZATION
%=====================================================================
\begin{column}{0.31\textwidth}
\seclabel{accentA}{1.\ MBQC Resource Optimization}
\hrule height 0pt

\begin{block}{Research Question}
  \RaggedRight
  How can a gate-based quantum circuit be translated into
  measurement-based quantum computation (MBQC) while
  \textbf{minimizing the required graph-state resources}?
\end{block}

\begin{block}{Motivation}
  \begin{itemize}
    \item Direct gate$\to$MBQC mapping produces large cluster states.
    \item Qubit count / entanglement is the practical bottleneck.
    \item Goal: an \emph{implementation-friendly} MBQC representation.
  \end{itemize}
\end{block}

\begin{block}{Method / Framework}
  \begin{itemize}
    \item Translate variational gate circuits into graph states.
    \item Optimizations currently explored:
      \begin{itemize}
        \item X-measurement chain compression
        \item removing redundant preparation entangling gates
        \item simpler single-qubit rotation decomposition
      \end{itemize}
    \item Treat measurement angles as variational parameters.
  \end{itemize}
\end{block}

\begin{block}{Current Direction / Insight}
  \begin{itemize}
    \item Preliminary reductions suggest parts of the graph
          resource are \textbf{redundant and compressible}.
    \item Structure of the source circuit guides where
          qubits can be saved.
  \end{itemize}
  \footnotesize\itshape (ongoing optimization -- results are preliminary)
\end{block}

\begin{block}{Next Direction}
  \begin{itemize}
    \item Systematic resource-reduction rules.
    \item Fully variational MBQC via measurement angles.
  \end{itemize}
\end{block}

\begin{block}{Key Schematic}
  \figplaceholder{accentA}{gate-based circuit $\to$ initial MBQC graph
    $\to$ optimized MBQC graph}
  \vskip1ex
  {\footnotesize\textbf{Fig.~1.} From a gate-based circuit to a reduced
   graph-state resource after successive compression steps.}
  % REPLACE with:
  % \includegraphics[width=0.95\linewidth]{figures/mbqc.pdf}
\end{block}
\end{column}

%=====================================================================
%  COLUMN 2 -- SURFACE-CODE SYNDROME UTILIZATION
%=====================================================================
\begin{column}{0.31\textwidth}
\seclabel{accentB}{2.\ Surface-Code Syndrome Utilization}
\hrule height 0pt

\begin{block}{Research Question}
  \RaggedRight
  Can \textbf{temporal syndrome information} identify dominant
  correlated faults that are \textbf{invisible in a single round}?
\end{block}

\begin{block}{Motivation}
  \begin{itemize}
    \item Distance-3 rotated surface code.
    \item In a single round, different CNOT fault locations can
          produce the \textbf{same syndrome signature}.
    \item This creates an information-theoretic ambiguity.
  \end{itemize}
\end{block}

\begin{block}{Method / Framework}
  \begin{itemize}
    \item Use \textbf{multi-round} detection-event sequences.
    \item Approaches explored:
      \begin{itemize}
        \item direct statistical separability analysis
        \item decoder-transformed features
        \item sequence models: RNN, GRU, Transformer
      \end{itemize}
  \end{itemize}
\end{block}

\begin{block}{Current Insight}
  \begin{itemize}
    \item Faults indistinguishable in a \textbf{single shot} become
          \emph{partially separable} once the \textbf{temporal
          sequence} is used.
    \item Time structure carries fault information the static
          syndrome hides.
  \end{itemize}
  \footnotesize\itshape (ongoing study -- qualitative, no fixed numbers yet)
\end{block}

\begin{block}{Next Direction}
  \begin{itemize}
    \item Quantify how much correlation is recoverable vs.\ distance/rounds.
    \item Link separability to decoder performance.
  \end{itemize}
\end{block}

\begin{block}{Key Schematic}
  \figplaceholder{accentB}{surface-code layout + syndrome-sequence
    heatmap; single-shot ambiguity vs.\ temporal separation}
  \vskip1ex
  {\footnotesize\textbf{Fig.~2.} Repeated syndromes over time separate
   faults that share one signature in a single round.}
  % REPLACE with:
  % \includegraphics[width=0.95\linewidth]{figures/surface.pdf}
\end{block}
\end{column}

%=====================================================================
%  COLUMN 3 -- VIRTUAL DISTILLATION THRESHOLD
%=====================================================================
\begin{column}{0.31\textwidth}
\seclabel{accentC}{3.\ Operational Threshold of Virtual Distillation}
\hrule height 0pt

\begin{block}{Research Question}
  \RaggedRight
  When does virtual distillation remain beneficial once the
  \textbf{purification circuit itself is noisy}?
\end{block}

\begin{block}{Motivation}
  \begin{itemize}
    \item Purification is usually analysed as ideal.
    \item Real purification circuits add their own gate noise.
    \item Core message: \textbf{purification must be reliable
          enough to stay useful.}
  \end{itemize}
\end{block}

\begin{block}{Method / Framework}
  \begin{itemize}
    \item Case study: noisy Bell-state purification.
    \item Include operational noise of the controlled-SWAP
          (Fredkin) gate.
    \item Also analyse Fredkin $\to$ CNOT decomposition with
          noisy CNOTs.
    \item Compare: ideal / noisy-coherent / destructive
          (ancilla-free) schemes.
  \end{itemize}
\end{block}

\begin{block}{Operational Threshold}
  \begin{itemize}
    \item Defined as the gate-noise level below which
          output fidelity is \textbf{actually improved}.
    \item Different decompositions / destructive schemes can show
          \textbf{qualitatively different} threshold behaviour.
  \end{itemize}
  \footnotesize\itshape (analytical + numerical work in progress;
  specific thresholds omitted here)
\end{block}

\begin{block}{Next Direction}
  \begin{itemize}
    \item Map thresholds across schemes and decompositions.
    \item Identify most noise-robust purification strategy.
  \end{itemize}
\end{block}

\begin{block}{Key Schematic}
  \figplaceholder{accentC}{output fidelity vs.\ operational gate-error
    rate, with the beneficial region marked by a threshold}
  \vskip1ex
  {\footnotesize\textbf{Fig.~3.} A threshold separates the regime where
   purification helps from where its own noise dominates.}
  % REPLACE with:
  % \includegraphics[width=0.95\linewidth]{figures/distillation.pdf}
\end{block}
\end{column}

\end{columns}

% ---- Unifying conclusion ban -----------------------------------------
\vskip0.4ex
\begin{beamercolorbox}[wd=\textwidth,sep=1.1ex,rounded=true,center]{block title}
  \normalsize
  \textbf{One theme.}\ \
  \textcolor{white}{Exploiting structure in quantum computation:}\ \
  \textcolor{white}{$\;$compress it (MBQC)$\;\;\longrightarrow\;\;$read it (syndromes)$\;\;\longrightarrow\;\;$restore it (distillation).}\ \
  Together, resource optimization, error inference, and noisy-state
  purification point toward more \textbf{efficient} and \textbf{reliable}
  quantum computation.
\end{beamercolorbox}

\end{frame}
\end{document}
